\newtoggle{withbonus} \togglefalse{withbonus}

\newcommand{\mycourse}{Ergänzungsfach Informatik}
\newcommand{\myclass}{4. Klassen}
\newcommand{\mydate}{12. November 2026}

\title{\textbf{Kolmogorov-Komplexität}}
\author{Thomas Graf\\
	\mycourse, \myclass}
\date{\mydate}
\maketitle
\thispagestyle{headandfoot}

\begin{center}
	\fbox{\parbox{140mm}{
		\begin{itemize}
			\item Dauer der Prüfung: $75$ Minuten
			\item Erlaubte Hilfsmittel: Schreibzeug, Taschenrechner
			\item Begründe deine Antworten und zeige deine Rechenschritte.
			\item Schreibe deine Lösungen auf separate Blätter.
		\end{itemize}}}
\end{center}
\vspace{10mm}

\begin{center}
	Vorname: \rule{45mm}{0.8pt}\qquad Nachname: \rule{45mm}{0.8pt}
\end{center}
\vspace{20mm}

\begin{center}
	\iftoggle{withbonus}{
		\combinedgradetable[h][questions]
	}{
		\gradetable[h][questions]
	}
\end{center}
\vspace{15mm}
\begin{center}
	Note:
	\vspace{10mm}

	\rule{8mm}{1.2pt}
\end{center}
\clearpage

% Lösungen werden über die Klassenoption answers in main.tex eingeblendet.
% Insgesamt: 22.5 Punkte.
\begin{questions}

	\question[2]
	Definiere ein neues \LaTeX-Makro mit dem Namen
	\lstinline[style=LaTeX]|begriff| und zwei Argumenten. Der Aufruf
	\lstinline[style=LaTeX]|\begriff{Komplexität}{Länge einer Beschreibung}|
	soll im normalen Fliesstext Folgendes erzeugen:
	\begin{center}
		\textbf{Komplexität}: \textit{Länge einer Beschreibung}
	\end{center}
	Das erste Argument soll fett, das zweite kursiv erscheinen. Der Doppelpunkt gehört zum Makro und soll nicht fett sein.
	\begin{solution}
		\begin{lstlisting}[style=LaTeX]
\newcommand{\begriff}[2]{\textbf{#1}: \textit{#2}}
		\end{lstlisting}
		\textbf{Bewertung:} 1 Punkt für eine korrekte Makrodefinition mit zwei Argumenten; 1 Punkt für Formatierung, Doppelpunkt und Abstand.
	\end{solution}

	\question[2]
	Das folgende vollständige \LaTeX-Dokument enthält vier voneinander unabhängige Fehler. Korrigiere den Code so, dass eine unnummerierte Überschrift und eine nummerierte Liste mit zwei Einträgen entstehen. Im zweiten Eintrag soll die Formel $x^2+1$ stehen.
	\begin{lstlisting}[style=LaTeX]
\usepackage{amsmath}
\documentclass{article}
\begin{document}
\section*{Unser Lernplan}
\begin{enumerate}
    \item Begriffe wiederholen.
    \item Die Formel x^2+1 untersuchen.
\end{itemize}
\end{article}
	\end{lstlisting}
	\begin{solution}
		\begin{lstlisting}[style=LaTeX]
\documentclass{article}
\usepackage{amsmath}
\begin{document}
\section*{Unser Lernplan}
\begin{enumerate}
    \item Begriffe wiederholen.
    \item Die Formel $x^2+1$ untersuchen.
\end{enumerate}
\end{document}
		\end{lstlisting}
		\textbf{Bewertung:} Je $0.5$ Punkte für die folgenden Korrekturen:
		\begin{itemize}
			\item Die Dokumentklasse steht vor dem Laden der Pakete.
			\item Die Formel steht im Mathematikmodus.
			\item Die Umgebung \texttt{enumerate} wird mit demselben Namen geschlossen.
			\item Das Dokument endet mit \lstinline[style=LaTeX]|\end{document}|.
		\end{itemize}
	\end{solution}

	\question
	\begin{parts}
		\part[1] Definiere die plain Kolmogorov-Komplexität $C(w)$ eines binären Wortes $w$. Gib an, in welcher Einheit sie gemessen wird.
		\begin{solution}
			$C(w)$ ist die minimale binäre Länge eines Programms, das ohne Eingabe genau $w$ ausgibt und anhält. Dabei ist eine optimale universelle Beschreibungssprache fest gewählt; im Skript verwenden wir die mit \texttt{g'++} kompilierten \texttt{C'++}-Programme. Gemessen wird die Länge des Maschinencodes in Bits.
		\end{solution}

		\part[2] Ein Programm mit einer binären Länge von $120$ Bits erzeugt ein Wort $w$ mit $10^6$ Bits und hält nach mehreren Tagen an. Eine Schülerin folgert: \enquote{Damit ist $C(w)=120$. Wegen der langen Laufzeit ist das Wort trotzdem nur schlecht komprimierbar.} Beurteile beide Aussagen.
		\begin{solution}
			Das Programm beweist nur $C(w)\leq120$. Ein noch kürzeres Programm könnte dasselbe Wort erzeugen; die Gleichheit ist deshalb nicht nachgewiesen. (1 Punkt)

			Die Laufzeit geht nicht in die Kolmogorov-Komplexität ein. Eine Beschreibung mit $120$ Bits ist gegenüber $10^6$ Bits sehr kurz, also ist $w$ im Sinne von Kolmogorov stark komprimierbar. (1 Punkt)
		\end{solution}
		\droptotalpoints
	\end{parts}

	\question
	Berechne die folgenden Summen. Gib die Ergebnisse als Ausdrücke ohne Summenzeichen an. Potenzen musst du nicht ausrechnen.
	\begin{parts}
		\part[1]
		\begin{align*}
			\sum_{k=0}^{30}\lr{3\cdot2^k}.
		\end{align*}
		\begin{solution}
			Mit der Formel für die geometrische Summe ergibt sich
			\begin{align*}
				\sum_{k=0}^{30}\lr{3\cdot2^k}
				=3\cdot\frac{2^{31}-1}{2-1}
				=3\lr{2^{31}-1}.
			\end{align*}
			\textbf{Bewertung:} $0.5$ Punkte für einen korrekten Rechenansatz und $0.5$ Punkte für das richtige Ergebnis.
		\end{solution}

		\part[1]
		\begin{align*}
			\sum_{k=1}^{20}4^k.
		\end{align*}
		\begin{solution}
			Die Summe beginnt bei $k=1$. Mit der Formel für die geometrische Summe folgt
			\begin{align*}
				\sum_{k=1}^{20}4^k
				=\frac{4^{21}-4^1}{4-1}
				=\frac{4^{21}-4}{3}.
			\end{align*}
			\textbf{Bewertung:} $0.5$ Punkte für einen korrekten Rechenansatz und $0.5$ Punkte für das richtige Ergebnis.
		\end{solution}
		\droptotalpoints
	\end{parts}


	\question[2]
	Bestimme die Längen der kürzesten Darstellung der beiden Zahlen $390624$ und $390625$ in Basis $5$.

	\textit{Hinweis:} $5^7=78125$ und $5^8=390625$.
	\begin{solution}
		Eine positive Zahl $m$ benötigt in Basis $b$ genau $k$ Ziffern, wenn $b^{k-1}\leq m<b^k$ gilt. Somit:
		\begin{align*}
			5^7\leq390624<5^8 &\quad\Longrightarrow\quad 8\text{ Ziffern},\\
			5^8\leq390625<5^9 &\quad\Longrightarrow\quad 9\text{ Ziffern}.
		\end{align*}
		Die erste Zahl hat die Darstellung $(44444444)_5$, die zweite $(100000000)_5$. Beim Übergang zur Potenz $5^8$ wird eine zusätzliche Stelle benötigt. Je 1 Punkt pro korrekt begründeter Länge.
	\end{solution}

	\clearpage
	\question[3]
	Für jedes natürliche $n\geq2$ sei $w_n=(001)^{n^2}$. Die folgenden Programmschemata sollen $w_n$ erzeugen. Welche sind gültig in \texttt{C'++}? Welche obere Schranke für $C(w_n)$ liefert jedes gültige Schema unmittelbar aus seiner Programmlänge, und welche dieser Schranken wächst am langsamsten mit $n$?

	\textbf{Vereinbarungen:} \texttt{Natural} bezeichnet beliebig grosse natürliche Zahlen mit exakter Arithmetik. Zwischen \texttt{***} und \texttt{+++} steht die konkrete Binärdarstellung von $n$ beziehungsweise $n^2$. Ein gültiges Programm darf nur einen verschiedenen binären Rohdatenblock enthalten; mehrfache Vorkommen desselben Blocks werden nur einmal gespeichert. Die Einbindung von \texttt{iostream} ist in allen drei Schemata vorausgesetzt.
	\begin{lstlisting}[language=C++,basicstyle=\ttfamily\small]
void An() {
    Natural n = 0b***n+++;
    for (Natural i = 0; i < n; ++i) {
        for (Natural j = 0; j < n; ++j) {
            std::cout << "001";
        }
    }
}

void Bn() {
    Natural m = 0b***n**2+++;
    for (Natural j = 0; j < m; ++j) {
        std::cout << "001";
    }
}

void Cn() {
    Natural n = 0b***n+++;
    Natural m = 0b***n**2+++;
    for (Natural j = 0; j < m; ++j) {
        std::cout << "001";
    }
}
	\end{lstlisting}
	\begin{solution}
		$A_n$ ist gültig und gibt den Block $001$ genau $n\cdot n=n^2$ Mal aus. Nur die Binärdarstellung von $n$ hängt von $n$ ab. Daher gilt
		\begin{align*}
			C(w_n)\leq\ceil{\log_2(n+1)}+c_A\leq\log_2(n)+c'_A.
		\end{align*}
		$B_n$ ist ebenfalls gültig und erzeugt dasselbe Wort. Hier wird jedoch $n^2$ gespeichert. Die unmittelbar erhaltene Schranke lautet
		\begin{align*}
			C(w_n)\leq\ceil{\log_2(n^2+1)}+c_B\leq2\log_2(n)+c'_B.
		\end{align*}
		Die Konstanten sind unabhängig von $n$. Die aus $A_n$ gewonnene Schranke hat den kleineren führenden Faktor; beide Schranken sind von der Grössenordnung $O(\log n)$.

		$C_n$ ist für $n\geq2$ ungültig: Die Binärdarstellungen von $n$ und $n^2$ sind verschieden. Dass die Variable \texttt{n} anschliessend nicht verwendet wird, ändert den Verstoss gegen die angegebene Sprachregel nicht.

		\textbf{Bewertung:} Je 1 Punkt für die Beurteilung von $A_n$ samt Schranke, von $B_n$ samt Schranke und Vergleich sowie von $C_n$ samt Begründung.
	\end{solution}

	\question[3]
	Für jedes natürliche $n\geq1$ sei $u_n=(10)^{2^{n^2}}$. Gib einen Algorithmus an, der nur $n$ als veränderlichen Datenblock speichert und ohne Eingabe $u_n$ erzeugt. Leite daraus eine obere Schranke für $C(u_n)$ her und drücke diese ausschliesslich durch die Wortlänge $L=\abs{u_n}$ und eine additive Konstante aus.

	Pseudocode genügt. Alle Berechnungen dürfen mit beliebig grossen natürlichen Zahlen exakt durchgeführt werden.
	\begin{solution}
		\begin{lstlisting}[language=C++,basicstyle=\ttfamily\small]
#include <iostream>

void Un() {
    Natural n = 0b***n+++;
    Natural exponent = n * n;
    Natural repetitions = 1;
    for (Natural j = 0; j < exponent; ++j) {
        repetitions = 2 * repetitions;
    }
    for (Natural j = 0; j < repetitions; ++j) {
        std::cout << "10";
    }
}
		\end{lstlisting}
		Nach der ersten Schleife ist \texttt{repetitions} gleich $2^{n^2}$. Die zweite Schleife erzeugt deshalb genau $u_n$. (1 Punkt)

		Der feste Programmteil ist unabhängig von $n$. Es folgt
		\begin{align*}
			C(u_n)\leq\ceil{\log_2(n+1)}+c_0\leq\log_2(n)+c_1.
		\end{align*}
		Dies ist eine obere Schranke, keine Behauptung über die exakte Komplexität. (1 Punkt)

		Da jeder Block zwei Bits enthält, gilt
		\begin{align*}
			L=2\cdot2^{n^2}=2^{n^2+1},
			\qquad n=\sqrt{\log_2(L)-1}.
		\end{align*}
		Einsetzen liefert die gesuchte Schranke
		\begin{align*}
			C(u_n)\leq\frac{1}{2}\log_2\lr{\log_2(L)-1}+c_1.
		\end{align*}
		Auch die etwas gröbere Schranke $C(u_n)\leq\frac12\log_2(\log_2 L)+c_1$ ist korrekt. (1 Punkt)
	\end{solution}

	\clearpage
	\question[1\half]
	Zeige: Für jedes natürliche $n\geq1$ haben mehr als sieben Achtel aller binären Wörter der Länge $n+3$ eine Kolmogorov-Komplexität von mindestens $n$.

	\textit{Hinweis:} Zähle alle möglichen binären Programme mit weniger als $n$ Bits. Berücksichtige dabei auch die Länge $0$.
	\begin{solution}
		Es gibt genau $2^{n+3}$ binäre Wörter der Länge $n+3$. Die Anzahl aller binären Zeichenfolgen, die als Programme mit weniger als $n$ Bits infrage kommen, beträgt
		\begin{align*}
			\sum_{k=0}^{n-1}2^k=2^n-1.
		\end{align*}
		($0.5$ Punkte)

		Jedes Programm erzeugt höchstens ein Wort. Manche Zeichenfolgen sind keine gültigen Programme, manche Programme halten nicht und manche erzeugen dasselbe Wort. Somit haben höchstens $2^n-1$ Wörter der betrachteten Länge eine Komplexität kleiner als $n$. ($0.5$ Punkte)

		Für den Anteil der Wörter mit $C(w)<n$ folgt daher
		\begin{align*}
			\frac{\#\{w\in\{0,1\}^{n+3}:C(w)<n\}}{2^{n+3}}
			&\leq\frac{2^n-1}{2^{n+3}}\\
			&=\frac18-\frac{1}{2^{n+3}}<\frac18.
		\end{align*}
		Weniger als ein Achtel der Wörter hat also eine Komplexität kleiner als $n$. Folglich haben mehr als sieben Achtel eine Komplexität von mindestens $n$. ($0.5$ Punkte)
	\end{solution}

	\question[3]
	Angenommen, ein Algorithmus $Q$ berechnet für jedes binäre Wort $w$ den exakten Wert $C(w)$ und hält immer an. Für $n\geq1$ betrachten wir folgenden Suchalgorithmus:
	\begin{enumerate}
		\item Durchlaufe alle binären Wörter der Länge $n$ in kanonischer Ordnung (mit $0<1$).
		\item Berechne für jedes Wort $w$ den Wert $Q(w)$.
		\item Sobald $Q(w)\geq n$ gilt, gib $w$ aus und halte an.
	\end{enumerate}
	Nenne das ausgegebene Wort $z_n$. Begründe, warum die Suche erfolgreich endet und weshalb die Beschreibung dieses Algorithmus zu einem Widerspruch führt. Den Satz über die Existenz inkompressibler Wörter darfst du voraussetzen.
	\begin{solution}
		Da nur Wörter derselben Länge $n$ betrachtet werden, stimmt die kanonische Ordnung hier mit der lexikographischen Reihenfolge überein. Für jede Länge $n$ existiert mindestens ein Wort $w$ mit $C(w)\geq n$. Die Suche prüft endlich viele Wörter, und jeder Aufruf von $Q$ hält nach Voraussetzung an. Deshalb wird ein solches Wort gefunden. Für das Ergebnis gilt $C(z_n)\geq n$. (1 Punkt)

		Die Aufzählung in kanonischer Ordnung, der Algorithmus $Q$ und die Ausgaberegel bilden einen festen Programmteil. Nur die Binärdarstellung von $n$ muss zusätzlich gespeichert werden. Daher gäbe es eine von $n$ unabhängige Konstante $c$ mit
		\begin{align*}
			C(z_n)\leq\ceil{\log_2(n+1)}+c.
		\end{align*}
		(1 Punkt)

		Zusammen müsste für jedes $n\geq1$ gelten:
		\begin{align*}
			n\leq C(z_n)\leq\ceil{\log_2(n+1)}+c.
		\end{align*}
		Für hinreichend grosse $n$ ist dies unmöglich, weil $n$ schneller wächst als $\log_2(n)$. Folglich kann ein solcher Algorithmus $Q$ nicht existieren. (1 Punkt)
	\end{solution}



	\question[1]
	Zeige, dass es für jede natürliche Zahl $n\geq1$ und jede ganze Zahl $i$ mit $0\leq i<n$ im Intervall $\ic{2^n}{2^{n+1}-1}$ mindestens $2^n-2^{n-i}$ natürliche Zahlen $x$ gibt mit $C(x)\geq n-i$. Dabei bezeichnet $C(x)$ die Kolmogorov-Komplexität der kürzesten Binärdarstellung von $x$.
	\begin{solution}
		Das Intervall $\ic{2^n}{2^{n+1}-1}$ enthält genau
		\begin{align*}
			(2^{n+1}-1)-2^n+1=2^n
		\end{align*}
		natürliche Zahlen. Jede davon hat eine kürzeste Binärdarstellung mit $n+1$ Bits.

		Da $n-i\geq1$ gilt, ist die Anzahl aller binären Zeichenfolgen der Länge kleiner als $n-i$ gleich
		\begin{align*}
			\sum_{k=0}^{n-i-1}2^k=2^{n-i}-1.
		\end{align*}
		Es gibt also höchstens $2^{n-i}-1$ Maschinencodes dieser Längen, die eine Ausgabe erzeugen. Jeder davon generiert höchstens eine Binärdarstellung und damit höchstens eine natürliche Zahl. Somit haben höchstens $2^{n-i}-1$ Zahlen im betrachteten Intervall eine Kolmogorov-Komplexität kleiner als $n-i$.

		Für die Anzahl $m$ der übrigen Zahlen folgt
		\begin{align*}
			m&\geq2^n-(2^{n-i}-1)\\
			 &=2^n-2^{n-i}+1\\
			 &>2^n-2^{n-i}.
		\end{align*}
		Dies beweist sogar die etwas stärkere untere Schranke $2^n-2^{n-i}+1$.

		\textbf{Bewertung:} $0.5$ Punkte für die begründete obere Schranke für die Anzahl der Zahlen mit $C(x)<n-i$ und $0.5$ Punkte für die korrekte Abschätzung von $m$.
	\end{solution}

\end{questions}
